\section{Morphisms of adelic curves}

We follow \cite[\S3.7]{chen-moriwaki}. Throughout this section,
$L^1(\Omega,\mathcal A,\nu)$ is the real vector space of integrable
real functions modulo equality almost everywhere. A morphism acts
contravariantly on fields and covariantly on places and fibre integrals.

\begin{definition}[Morphism of adelic curves]
\label{def:adelic-morphism}
\leanok
\source{chen-moriwaki-adelic:page-0219}
\uses{def:adelic-structure}
\lean{ArakelovGeometry.AdelicCurve.Hom,ArakelovGeometry.AdelicCurve.Hom.pullback,ArakelovGeometry.AdelicCurve.Hom.coeFn_pullback,ArakelovGeometry.AdelicCurve.Hom.integrable_comp,ArakelovGeometry.AdelicCurve.Hom.integral_comp,ArakelovGeometry.AdelicCurve.Hom.ext,ArakelovGeometry.AdelicCurve.Hom.measurable_baseMap,ArakelovGeometry.AdelicCurve.Hom.map_measure,ArakelovGeometry.AdelicCurve.Hom.fibreIntegral_pullback_apply}
Let $S'=(K',(\Omega',\mathcal A',\nu'),\phi')$ and
$S=(K,(\Omega,\mathcal A,\nu),\phi)$ be adelic curves.
A morphism $\alpha:S'\to S$ is a triple
$(\alpha^\#,\alpha_\#,I_\alpha)$ with the following properties:
\begin{enumerate}
\item $\alpha^\#:K\to K'$ is a field homomorphism.
\item $\alpha_\#: \Omega'\to\Omega$ is measurable,
$(\alpha_\#)_*\nu'=\nu$, and, for every $\omega'\in\Omega'$ and
$a\in K$,
\[
 \phi'_{\omega'}(\alpha^\#(a))
   =\phi_{\alpha_\#(\omega')}(a).
\]
\item $I_\alpha:L^1(\Omega',\mathcal A',\nu')
\to L^1(\Omega,\mathcal A,\nu)$ is real-linear and satisfies
\[
 \int_\Omega I_\alpha(g)\,d\nu
   =\int_{\Omega'}g\,d\nu'
 \quad(g\in L^1(\Omega',\mathcal A',\nu')),
 \qquad
 I_\alpha\alpha^*=\operatorname{id}_{L^1(\Omega,\mathcal A,\nu)},
\]
where $\alpha^*[f]=[f\circ\alpha_\#]$ denotes pullback.
\end{enumerate}
In particular, for every integrable real $f$ on $\Omega$,
$f\circ\alpha_\#$ is integrable and
\[
 \int_{\Omega'}f\circ\alpha_\#\,d\nu'
   =\int_\Omega f\,d\nu.
\]
Two morphisms are equal precisely when their three data are equal.
Only real linearity is required of $I_\alpha$; positivity and
continuity are not additional morphism axioms.
\end{definition}
\begin{proof}
The pullback appearing in the definition is well-defined as follows.
For a measurable integrable $f$, the pushforward equality gives
$\int_{\Omega'}|f\circ\alpha_\#|\,d\nu'
=\int_\Omega|f|\,d\nu<\infty$.
If $f_1=f_2$ almost everywhere, the exceptional set is contained in a
measurable null set $N$. Its inverse image is measurable and has measure
$\nu(N)=0$, so the two pullbacks agree almost everywhere.
For an arbitrary integrable representative, choose a measurable
representative and apply this null-set argument. This proves both
integrability and independence of representatives. Pullback preserves
addition and scalar multiplication pointwise. Applying the pushforward
formula to the positive and negative parts proves the displayed equality
of real integrals. Thus the class identity in (3) says exactly that
$I_\alpha$ sends the class of $f\circ\alpha_\#$ to the class of $f$.
\end{proof}

\begin{lemma}[Identity morphism]
\label{lem:adelic-morphism-id}
\leanok
\source{chen-moriwaki-adelic:page-0219}
\uses{def:adelic-morphism}
\lean{ArakelovGeometry.AdelicCurve.Hom.id}
For every adelic curve $S=(K,(\Omega,\mathcal A,\nu),\phi)$,
the triple
\[
 \operatorname{id}_S=
 (\operatorname{id}_K,\operatorname{id}_\Omega,
  \operatorname{id}_{L^1(\Omega,\mathcal A,\nu)})
\]
is a morphism from $S$ to itself.
\end{lemma}
\begin{proof}
The identity of $K$ is a field homomorphism. The identity of $\Omega$
is measurable and preserves $\nu$. Its compatibility with the absolute
values is $\phi_\omega(a)=\phi_\omega(a)$.
The identity of $L^1$ is real-linear and preserves every integral.
Pullback by $\operatorname{id}_\Omega$ fixes each representative, so
the identity of $L^1$ is also a left inverse to that pullback.
\end{proof}

\begin{lemma}[Composition of morphisms]
\label{lem:adelic-morphism-comp}
\leanok
\source{chen-moriwaki-adelic:page-0219}
\uses{def:adelic-morphism}
\lean{ArakelovGeometry.AdelicCurve.Hom.comp,ArakelovGeometry.AdelicCurve.Hom.pullback_comp}
Let $\alpha:S'\to S$ and $\beta:S''\to S'$ be morphisms of
adelic curves. Then
\[
 \alpha\circ\beta=
 (\beta^\#\circ\alpha^\#,
  \alpha_\#\circ\beta_\#,
  I_\alpha\circ I_\beta)
\]
is a morphism from $S''$ to $S$.
\end{lemma}
\begin{proof}
Composition preserves field homomorphisms and measurable maps.
Writing $\nu'',\nu',\nu$ for the three measures, functoriality of
pushforward gives
\[
 (\alpha_\#\circ\beta_\#)_*\nu''
 = (\alpha_\#)_*((\beta_\#)_*\nu'')
 = (\alpha_\#)_*\nu'=\nu.
\]
For $\omega''\in\Omega''$ and $a\in K$, the two absolute-value
compatibilities give
\[
 \phi''_{\omega''}(\beta^\#(\alpha^\#(a)))
 =\phi'_{\beta_\#(\omega'')}(\alpha^\#(a))
 =\phi_{\alpha_\#(\beta_\#(\omega''))}(a).
\]
The composite $I_\alpha I_\beta$ is real-linear. For every
$g\in L^1(\Omega'',\mathcal A'',\nu'')$, its integral is
\[
 \int_\Omega I_\alpha(I_\beta g)\,d\nu
 =\int_{\Omega'}I_\beta g\,d\nu'
 =\int_{\Omega''}g\,d\nu''.
\]
On representatives,
$(\alpha\circ\beta)^*=\beta^*\alpha^*$.
Consequently, for $f\in L^1(\Omega,\mathcal A,\nu)$,
\[
 I_\alpha I_\beta(\beta^*\alpha^*f)
 =I_\alpha(\alpha^*f)=f.
\]
All the morphism axioms follow.
\end{proof}

\begin{theorem}[The category of adelic curves]
\label{thm:adelic-curve-category}
\leanok
\source{chen-moriwaki-adelic:page-0219}
\uses{def:adelic-morphism,lem:adelic-morphism-id,lem:adelic-morphism-comp}
\lean{ArakelovGeometry.AdelicCurve.Hom.id_comp,ArakelovGeometry.AdelicCurve.Hom.comp_id,ArakelovGeometry.AdelicCurve.Hom.comp_assoc,ArakelovGeometry.AdelicCurve.category}
Adelic curves, with the morphisms of Definition~\ref{def:adelic-morphism},
the identities above, and the stated composition, form a category.
\end{theorem}
\begin{proof}
The preceding lemmas supply the identity at each object and composition
for every composable pair. A morphism is determined by its field map,
place map, and linear map. Composing each of these three maps with its
identity leaves it unchanged, on either side. This proves the two unit
laws. For morphisms $\gamma:S'''\to S''$, $\beta:S''\to S'$, and
$\alpha:S'\to S$, both bracketings have field map
$\gamma^\#\circ\beta^\#\circ\alpha^\#$, place map
$\alpha_\#\circ\beta_\#\circ\gamma_\#$, and linear map
$I_\alpha\circ I_\beta\circ I_\gamma$.
Associativity of composition of functions therefore proves the
associativity law for morphisms.
\end{proof}

\begin{example}[The morphism from an algebraic extension]
\label{ex:algebraic-extension-morphism}
\leanok
\source{chen-moriwaki-adelic:page-0220}
\uses{def:adelic-morphism,def:algebraic-extension,cor:algebraic-L1-operator,lem:algebraic-base-pushforward,prop:algebraic-signed-disintegration,lem:algebraic-integral-pullback,prop:algebraic-base-pullback}
\lean{ArakelovGeometry.AdelicCurve.Hom.algebraicExtension,ArakelovGeometry.AdelicCurve.Hom.algebraicExtension_fibreIntegral_toL1,ArakelovGeometry.AdelicStructure.Algebraic.measurePreserving_proj,ArakelovGeometry.AdelicStructure.Algebraic.integral_fibreIntegralL1,ArakelovGeometry.AdelicStructure.Algebraic.fibreIntegralL1_compMeasurePreserving,ArakelovGeometry.AdelicCurve.Hom.algebraicExtension_fieldMap,ArakelovGeometry.AdelicCurve.Hom.algebraicExtension_baseMap,ArakelovGeometry.AdelicCurve.Hom.algebraicExtension_fibreIntegral,ArakelovGeometry.AdelicCurve.Hom.coeFn_algebraicExtension_fibreIntegral}
Let $S=(K,(\Omega,\mathcal A,\nu),\phi)$ be an adelic curve and
let $L/K$ be an algebraic field extension. For its adelic extension
$S\otimes_K L$, the triple
\[
 (K\hookrightarrow L,\pi_{L/K},\overline I_{L/K})
 :S\otimes_K L\longrightarrow S
\]
is a morphism, where
$\overline I_{L/K}[f]=[I_{L/K}f]$ is the induced real-linear
operator on $L^1$.
\end{example}
\begin{proof}
The field inclusion is a homomorphism. The projection $\pi_{L/K}$
is measurable by the construction of the adelic extension, and
Lemma~\ref{lem:algebraic-base-pushforward} gives
$(\pi_{L/K})_*\nu_L=\nu$. A place of $S\otimes_K L$ over
$\omega$ is an absolute value on $L$ whose restriction to $K$ is
$\phi_\omega$. This is precisely the required absolute-value
compatibility.

Corollary~\ref{cor:algebraic-L1-operator} gives the real-linear
operator $\overline I_{L/K}$ with the asserted representative formula.
For every integrable $f$ on $\Omega_L$,
Proposition~\ref{prop:algebraic-signed-disintegration} gives
\[
 \int_\Omega\overline I_{L/K}[f]\,d\nu
 =\int_\Omega I_{L/K}f\,d\nu
 =\int_{\Omega_L}f\,d\nu_L.
\]
For integrable $h$ on $\Omega$, the pullback $h\circ\pi_{L/K}$
is integrable by Proposition~\ref{prop:algebraic-base-pullback}.
Lemma~\ref{lem:algebraic-integral-pullback} says
$I_{L/K}(h\circ\pi_{L/K})=h$. Passing to equivalence classes gives
$\overline I_{L/K}[h\circ\pi_{L/K}]=[h]$.
Thus both axioms for the fibre integration operator hold.
\end{proof}

\begin{example}[The morphism to a subfield restriction]
\label{ex:subfield-restriction-morphism}
\leanok
\source{chen-moriwaki-adelic:page-0220}
\uses{def:adelic-morphism,def:subfield-restriction}
\lean{ArakelovGeometry.AdelicCurve.Hom.fRestriction,ArakelovGeometry.AdelicCurve.Hom.fRestriction_fieldMap,ArakelovGeometry.AdelicCurve.Hom.fRestriction_baseMap,ArakelovGeometry.AdelicCurve.Hom.fRestriction_fibreIntegral}
Let $S=(K,(\Omega,\mathcal A,\nu),\phi)$ be an adelic curve and
let $K_0\subseteq K$ be a subfield. There is a morphism
\[
 (K_0\hookrightarrow K,\operatorname{id}_\Omega,
       \operatorname{id}_{L^1(\Omega,\mathcal A,\nu)})
 :S\longrightarrow S|_{K_0}.
\]
\end{example}
\begin{proof}
The inclusion of the subfield is a field homomorphism. The source and
target have the same measure space, so the identity on $\Omega$ is
measurable and preserves the measure. For $a\in K_0$, the absolute value
of its image in $K$ equals its restricted absolute value by the definition
of $S|_{K_0}$. The two real $L^1$ spaces are identical. Their identity
map is real-linear and preserves every integral, and pullback along the
identity place map is the identity. Hence it has the required left inverse.
\end{proof}
